\subsection{Native recovery with an online selected message}
\label{sec:native-online-chain}

The finite-chain comparison can now be connected to the original native
recovery experiment. The connection preserves its acquisition-based
eligibility and executed source slots while allowing earlier message-key
disclosures to influence which later application send is challenged.
The root, traffic material and message step have separate allocation
boundaries; keeping those boundaries explicit makes the selection
implementable by the same native execution ledger.

\paragraph{The online native controller.}
Retain the established-session live core of
Section~\ref{sec:live-core-instantiation}, with one authoritative history
per endpoint and passive copies, and fix $1\le L\le2^{64}-1$ before
interaction. For each fixed candidate $j=(g,a)$, designate its original
effective tuple $m^*$ immediately before its first private ContextBound
hash computation at any protected alias, as in
Section~\ref{sec:contextbound-connection}. Designate the complete
traffic entry $I^*$ and direction before that entry's first block allocation.
These designations are retained through later acceptance, recovery and
equal-input aliases.

The controller maintains a private bookkeeping table indexed by
$(\mathcal M,d,i)$, where $\mathcal M$ is the complete original traffic
material handle, $d$ its direction and $i$ its chain index. Equal-info
aliases of the same root schedule share $\mathcal M$; an independently
allocated material has its own handle even when its epoch, index or
numeric key bytes coincide. Immediately before a previously unseen
selected-direction step, the controller chooses either Continue or
Mark, using its available history and the public descriptor of that
step. Mark occurs at most once and only for $i<L$. A receive-side
advance, including a skipped step followed by failed authentication,
is a first computation for this purpose. Repeating a computation after
a failed write retrieves the existing ledger entry and offers no new
selection point. This ledger is reduction bookkeeping and changes no
native state transition or response.

For each step passed with Continue before Mark, the observer may obtain
its message key; the prefix seed and Extract value retain private
handles. At a marked
step $\ell$, the observer may obtain the paired successor seed
$s_{\ell+1}$ after its derivation. This named projection is bound to
$(\mathcal M,d,\ell)$, irrespective of subsequent progress at either
endpoint. Later chain steps are real computations from that seed.
The selected current root, private prefix seeds, Extract values and
marked message key have only the uses in the preceding interfaces.
The already admitted successor-root, opposite-direction and other
traffic-entry observations remain available.

These typed projections are the complete additional observation
interface. Acquiring a complete session version still has its original
scope and save time. In particular, a version containing an available
successor seed also contains an acquired session version and remains
subject to the original acquisition predicate. Public commitments and
failed captures retain their original behavior.

There is one left-or-right application operation, tied to the marked
step's original send. Its messages have equal length and common
associated data. A native rejection, uncertain outcome or exact retry
does not move that operation to another step. The terminal selector
chooses the cut and original reservation handles as before and must
identify this same accepted logical send. If no mark or accepted send
is reached, the original valid-target condition is false. Application
payloads and payload-bearing delivery remain private; the adversary
has the existing plaintext-free Resume and public protocol interfaces.
Thus $E$ is exactly the original eligibility predicate evaluated on
this controller's execution.

A mark may first occur during receive-side advancement, before the
target send is reached. Accordingly the one-key AEAD experiments below
allow classical decryption queries before their at-most-one encryption
call. Such earlier queries retain their actual replies. After the
left-or-right call, an exact challenge decryption returns the existing
authenticated marker. This operational convention preserves the
native receive-before-send order and is included in the stated AEAD
advantages.

\begin{lemma}[Native online-chain correspondence]
\label{lem:native-online-chain-correspondence}
For this controller and a fixed candidate $j$, the stopped native live
core realizes the prefix interface of Section~\ref{sec:chain-prefix}
after the existing record, component, ContextBound and root-entry
comparisons. Its original send and receive computations use one cached
step per $(\mathcal M,d,i)$. The marked step has at most one original
encryption tuple, and the live-core charged projection event remains
empty.
\end{lemma}
\begin{proof}
At traffic installation, the exact root schedule and complete traffic
information string determine the material handle. The role-dependent
assignment in the native traffic constructor makes one endpoint's
send seed the matching endpoint's receive seed; the opposite direction
has its other handle. Copying, reopening and exact pending-target
reconciliation preserve this association. The source-attribution
induction fixes the KEM origins of each installed schedule independently
of which message index is later marked.

Induct on actual private chain calls across both endpoints. For index
zero, the constructor supplies the corresponding traffic block. A
later call uses the predecessor's successor with the fixed native
domain and big-endian index. A repeated call under the same material,
direction and index therefore repeats the same Extract/Expand inputs.
Record its result at first computation, including in transaction-private
receive state. If receive authentication or a later durable update
fails, the native state may return to an earlier counter, while this
proof ledger retains the original result. A retried advance obtains
that result again. Taking an already retained skipped key likewise
introduces no new computation or marking opportunity. This is an
exact representation of the deterministic native computation in
each comparison world.

The sender stores its pending message plan before reaching the chain
step and encryption. A first-retention failure exits before encryption;
once encryption is reached, the authoritative plan binds the message
identifier, plaintext and associated data. A failed later write or
process loss recovers that same plan or its committed outbox. Exact
resumption uses its original encryption tuple, and successful commitment
advances the send index. Hence the first encryption answer can be
cached before its durable outcome, with at most one original tuple
under the marked material and index. This also covers a mark whose
first step computation occurred at the receiver.

The representation invariant for images and intents carries the same
material and step handles through their immutable versions. Acquiring
an intact version of an occurred candidate reservation's session,
saved at or after that reservation, certifies the existing free stop.
Lemma~\ref{lem:adaptive-origin-stop} applies before a final cut is chosen:
every eligible cut precedes both reservations and therefore that
acquired version's save. This certificate is independent of message
index or marking time. Complete old or unrelated versions
retain their ordinary correspondence. The new individual projections
return exactly earlier message blocks and the marked pair's successor,
which the prefix reduction already supplies. They expose no additional
protected record or report key. Thus the live core retains
$B_j=\varnothing$ and its established real-to-opaque record comparison.

After the ContextBound and root-entry comparisons, the complete native
continuation therefore has the three online designations and the
private prefix uses of Theorem~\ref{thm:online-chain-leaf}. Real control
verification, ACKs, other materials and later epochs remain in that
continuation. The whole-execution budgets include their computations
and all newly permitted prefix and suffix observations.
\end{proof}

\paragraph{Original-game costs.}
Use the original native $E,A,I$ and its residual
$\eta_{\rm attr}^{\rm nat}=\Pr[E\wedge\neg A]$. Let
$\varepsilon_{{\rm rec},j}$ be the record event cost in
\eqref{eq:live-core-record-cost}, now with the complete online controller's
budgets. The operational component adapters retain their definitions
from Section~\ref{sec:contextbound-connection}, with their output events
evaluated in this expanded experiment. For candidate $j$, set
\begin{equation}
 B_j^{\rm chain}=\min\bigl\{1,\;
  2\Phi_{256}(d_j,Q_{H,j})+2\Delta_{{\rm init},j}
  +2\Delta_{{\rm prefix},j}(L)+2\varepsilon_{{\rm AEAD},j}\bigr\}.
 \label{eq:native-chain-continuation}
\end{equation}
Every budget is fixed before interaction and covers preparation,
failed attempts, all endpoint aliases, the adaptive mark, real suffix
computation and the terminal test in its corresponding comparison.
The prefix sum in \eqref{eq:prefix-error} gives the explicit index-wise
comparison cost. No additional combinatorial message-index guessing
factor is introduced; all other whole-execution budgets may also grow
with $L$.

Let $\mathsf{Bad}_{\rm msg}^{\rm chain}$ record a fresh authenticated
receive and inbox installation under the marked step whose complete
$(\mathrm{AAD},\mathrm{ciphertext})$ pair is absent from the earlier
honest encryption ledger for that material, direction and index.
The flag is recorded at that installation, including before the target
send, and later encryption does not erase it. Accepted-send provenance
identifies the event at the end of the original experiment. For the
component adapter with Boolean output
$W_j=\mathbf1[E\wedge A\wedge I=j\wedge
\mathsf{Bad}_{\rm msg}^{\rm chain}]$, write its operational advantage
as $a_j^{\rm msg}$ and define
\begin{equation}
 \Omega_j^{\rm chain}=\min\bigl\{1,\;
  \Phi_{256}(d_j,Q_{H,j})+\Delta_{{\rm init},j}
  +\Delta_{{\rm prefix},j}(L)+\varepsilon_{{\rm INT},j}\bigr\}.
 \label{eq:native-chain-integrity-cost}
\end{equation}
The costs for the two terminal tests use their own complete-query
bounds; sharing notation does not identify their distinguishers.

\begin{theorem}[Native online selected-message security]
\label{thm:native-online-message}
In the online native live-core experiment above,
\begin{equation}
\begin{split}
 \operatorname{Adv}_{\rm msg}^{\rm chain}
 \le\min\biggl\{\frac12,\frac12\Bigl[
    \eta_{\rm attr}^{\rm nat}\\[-2pt]
 \qquad+
 \sum_{j\in J}\bigl(2\varepsilon_{{\rm rec},j}
       +a_j^++a_j^-+2\delta_j+B_j^{\rm chain}\bigr)
                       \Bigr]\biggr\},
\end{split}
\label{eq:native-online-confidentiality}
\end{equation}
and
\begin{equation}
\begin{split}
 \Pr_{G_{\rm nat}}[E\wedge\mathsf{Bad}_{\rm msg}^{\rm chain}]
 \le\min\biggl\{1,\;&\eta_{\rm attr}^{\rm nat}\\[-2pt]
 &+\sum_{j\in J}\min\{1,\varepsilon_{{\rm rec},j}
       +a_j^{\rm msg}+\delta_j+\Omega_j^{\rm chain}\}\biggr\}.
\end{split}
\label{eq:native-online-integrity}
\end{equation}
For the executed-slot selector, the corresponding original-game
residual in each bound is at most $\Delta_{\rm attr}$ from
\eqref{eq:executed-attribution-cost}, using the complete expanded
signing experiment and its established reservation-signature bridge.
\end{theorem}
\begin{proof}
For confidentiality, put
$S=2\mathbf1[g=b]-1$ and run the same available candidate terminal
code $\Psi_j=\mathbf1[E\wedge A\wedge I=j]S$ in every world.
The identity in Lemma~\ref{lem:native-local-score} localizes the
original score before any candidate-specific comparison. The native
free stops remain valid by
Lemma~\ref{lem:native-online-chain-correspondence}. Its record
correspondence gives $\rho_j\le2\varepsilon_{{\rm rec},j}$;
the two operational component-event comparisons cost
$a_j^++a_j^-+2\delta_j$.

In the random-component world, apply
Lemma~\ref{lem:cb-hidden-component} with the unchanged complete tuple
cache. The target contribution is selected at its own first private
allocation, before the later traffic and message designations. Its
single root-use entry is still that of
Lemma~\ref{lem:cb-native-root-use}, since later chain work does not
reuse the combined contribution or alter its original pending-root
inputs. The two score comparisons therefore cost
$2\Phi_{256}(d_j,Q_{H,j})+2\Delta_{{\rm init},j}$.
The resulting uniform-root continuation is exactly the online prefix
interface by the new correspondence lemma.
Corollary~\ref{cor:online-chain-message} bounds its qualified signed
score by $2\Delta_{{\rm prefix},j}(L)+2\varepsilon_{{\rm AEAD},j}$.
Every selector, disclosure, eligibility check and terminal event is
recomputed within that world. In particular, qualifying a run uses
the available provenance and view rather than direct access to $b$.
Thus the random-component score is bounded by $B_j^{\rm chain}$.
Substitution in Proposition~\ref{prop:native-local-accounting}
proves~\eqref{eq:native-online-confidentiality}.

For integrity, apply the Boolean localization identity from
\eqref{eq:native-auth-localization} to the selected-message event
$W_j$. The same free-stop certificate makes a stopped candidate's
final conjunction false, even if its recorded inbox event preceded
the challenge send. The record comparison costs
$\varepsilon_{{\rm rec},j}$, and the operational component comparison
costs $a_j^{\rm msg}+\delta_j$. Apply the complete-tuple and root-entry
comparisons to this Boolean continuation. The remaining event is a
subset of the selected-key integrity event in
Corollary~\ref{cor:online-chain-message}, so its probability is bounded
by $\Omega_j^{\rm chain}$. A fresh successful decryption before the
honest encryption is already a valid integrity witness; the reduction
may return it immediately. This justifies the same bound when Mark
was first reached by the receiver. Summing the native candidate
partition and capping by one proves~\eqref{eq:native-online-integrity}.

Finally, the executed-slot association proof is an induction over
traffic installation and its originating KEM plans. Its signed-body
and full-wire checks precede the message-index choice. Advancing,
retaining or reopening the installed traffic entry preserves those
origins, while the earlier-key projections only enlarge the complete
adversarial continuation. Proposition~\ref{prop:executed-source-attribution}
therefore gives the stated original-game attribution bound with
its updated whole-execution resources. This term is charged once,
and the root, record and component costs retain their displayed
coefficients; no additional guess of the marked message index is used.
\end{proof}

\paragraph{Reading the bound in ordinary primitives.}
For the executed-slot selector, Table~\ref{tab:native-online-substitutions}
collects the substitutions in
\eqref{eq:native-online-confidentiality}--\eqref{eq:native-online-integrity}.
The index $\sigma\in\{+,-,\mathrm{msg}\}$ identifies the actual terminal
test: each test retains its own complete-execution resources. The table
instantiates the displayed terms without changing the original candidate
partition or the selected message.

\begin{table}[htbp]
\caption{Primitive substitutions for native online selected-message security.
The attribution cost is charged once outside the candidate sum.}
\label{tab:native-online-substitutions}
\small
\begin{tabularx}{\textwidth}{@{}p{0.19\textwidth}>{\raggedright\arraybackslash}X@{}}
\toprule
Term & Established substitution and its role \\
\midrule
Record cost $\varepsilon_{{\rm rec},j}$ &
$\displaystyle\min\!\left\{1,
\sum_{i=1}^{M_j^{\rm seal}}\Phi_{192}(d_{j,i},Q_{j,i})
+\delta_{{\rm seal},j}\right\}$.
\par Here $\delta_{{\rm seal},j}=\min\{1,\binom{M_j^{\rm seal}}2 2^{-192}\}$.
Equation~\eqref{eq:live-core-record-cost} counts fresh sealing attempts,
including failed commits; copies retain their original handles. \\
\addlinespace[3pt]
Component cost $a_j^\sigma$ &
At most $\min\{1,\operatorname{Adv}^{\rm cca}(B_j^\sigma)
+\zeta_{{\rm pre},j}^{\sigma}\}$.
\par Here
$\zeta_{{\rm pre},j}^{\sigma}
=\min\{1,q_{j,\sigma}^{-}2^{-256}+\delta_{{\rm mr},j}^{\sigma}\}$
by Proposition~\ref{prop:kem-deferred-challenge}.
The bound retains real prechallenge decapsulation. \\
\addlinespace[3pt]
Correctness $\delta_j$ &
The selected pair's shared-key correctness cost from
Section~\ref{sec:contextbound-connection}. It remains separate from the
message-recovery event $\delta_{{\rm mr},j}^{\sigma}$ inside the component
bridge. \\
\addlinespace[3pt]
Continuation $B_j^{\rm chain}$,
$\Omega_j^{\rm chain}$ &
Equations~\eqref{eq:native-chain-continuation} and
\eqref{eq:native-chain-integrity-cost} combine ContextBound replacement,
root initialization, the finite prefix sum in~\eqref{eq:prefix-error},
and one-key AEAD confidentiality or integrity. Each comparison costs twice
its event bound for the signed score and once for the Boolean event. \\
\addlinespace[3pt]
Attribution $\Delta_{\rm attr}$ &
$\min\{1,\varepsilon_{\rm sig}+\chi_H(Q_H)\}$ by
Proposition~\ref{prop:executed-source-attribution}, using the ordinary
ML-DSA instantiation below. This is the original-game association cost
for the executed-slot selector. \\
\bottomrule
\end{tabularx}
\end{table}

Specifically, for at most $U$ fresh signing-owner families, the established
ordinary-signature bridge~\eqref{eq:native-pair-ordinary-mldsa} permits
\[
\begin{split}
 \varepsilon_{\rm sig}\le\min\bigl\{1,U\bigl(&
 \varepsilon_{\rm MLDSA,EUF}^{\rm su}+\beta_{256}(q)\\[-2pt]
 &+\delta_{\mathsf X}(q_*^{\rm res})
  +\delta_{\mathsf X}(q_*^{\rm ord})\bigr)\bigr\}.
\end{split}
\]
Here $q$ bounds distinct executed target $(r,\mu)$ pairs,
$\beta_{256}(q)=\min\{1,\binom q2 2^{-256}\}$, and the two $q_*$ budgets
count remaining shared-XOF calls in the complete reservation and ordinary
experiments. The established oracle terms are
$\delta_{\mathsf C}(v)=\min\{1,v2^{-256}\}$ in the classical model and
$\delta_{\mathsf Q}(v)=\min\{1,2v2^{-128}\}$ in the quantum model.

Confidentiality uses the two component tests $+$ and $-$, each with its
own prechallenge term, before the outer factor $1/2$. Integrity uses the
single $\mathrm{msg}$ test. The existing $2\delta_j$ and $\delta_j$
coefficients therefore remain in their respective bounds. The live core
has zero charged projection residual. Online marking incurs the displayed
$L$-term prefix sum; its complete-query budgets include preparation,
failed attempts, endpoint aliases, real suffix computation and terminal
testing. These are the same interfaces and budgets as in the theorem.

For $L=1$, the prefix cost reduces to the existing first-message
schedule cost. The enlarged controller adds the choice of a later
uncomputed index after observing earlier keys, while preserving the
native source and acquisition semantics. An explicitly modeled
record/report projection extension keeps its original
$\eta_{\rm proj}^{\rm nat}$ and corresponding local comparison costs,
as in the preceding native theorems.
